\chapter[Teufelsanbeter]{Teufelsanbeter}
\quellseite{8}
Daß es in unsern Tagen noch Teufelsanbeter oder sogenannte „Satanianer“ gibt, dafür lieferte neulich eine telegraphische Nachricht aus Paris den unumstößlichsten Beweis. In derselben wurde berichtet, daß die Polizei eine große Versammlung von Satanianern überrascht und ausgehoben habe, als sie eben die schändlichsten Orgien feierten. Ganz unglaubliche Sittlichkeitsverbrechen gehören nach den Aussagen dieser Teufelsanbeter zu ihrem Cultus- oder Teufelsdienst. Die Hauptfeier, die sie ihrem Herrn, dem Satan, veranstalten, nennen sie „schwarze Messe“, die mit furchtbaren Ceremonien beginnt und mit scheußlichen Unsittlichkeiten endet. Ein ungläubiger französischer Zeitungsschreiber, der die Feier der „schwarzen Messe“ in seinem Blatte geleugnet und als eine fixe Idee und grundlose Beschuldigung der Christen bezeichnet hatte, erhielt eines Abends die Einladung, in einem näher bezeichneten Hause um 10 Uhr zu erscheinen, um einer „schwarzen Messe“ beizuwohnen. Er leistete der Einladung Folge. Zur bestimmten Zeit trat er in ein nur schwach erleuchtetes Zimmer, dessen Wände mit dunkelrothem Tuch behangen waren. Es gewahrte allerlei phantastische, grauenerregende Figuren und eine aus Gyps verfertigte Christusstatue. Plötzlich wurde der Saal durch elektrisches Licht grell erleuchtet: vor einem rothbehangenem „Altar“ stand

\quellseite{9}
ein rothgekleideter Satanspriester, der die Augen furchtbar verdrehte, das Gesicht schrecklich verzerrte, die Hände krampfhaft faltete: er „betete“ zum Teufel. Der ganze Saal war dicht besetzt von Männern, Frauen und Kindern, die ihre Gesichter scheußlich geschminket und tätowirt hatten. Jetzt treten plötzlich fratzenhaft bemalte Chorknaben ein und stimmen den „Meßgesang“ an. Die ganze Versammlung stimmt ein: Anfangs leise, allmählich lauter. Der Priester verläßt den „Altar“ und mit grimmigen Geberden wendet er sich dem Christusbilde zu, erhebt mit der Linken sein Meßgewand, mit der Rechten ergreift er einen \unsicher{Hämmer} und nähert sich dem verhaßten Bilde, furchtbare Flüche wider Christum ausstoßend. Die ganze Versammlung tobet auf in höllischem Haß und stößt unaussprechliche Lästerungen und unflätige Schmutzreden wider Christum aus. Dem Zeitungsschreiber standen die Haare zu Berge, die Knie schlotterten, und der Schweiß drang ihm aus allen Poren. Er hätte sich gerne weggeschlichen, aber er war wie versteinert. Plötzlich trat Todtenstille ein: der Priester stand unmittelbar vor dem Christusbilde, er erhob den Hammer und unter gräßlichen Flüchen schlug er so lange auf die Statue ein, bis sie zertrümmert zu Boden stürzte. Die Menge brach in lautes Beifallsschreien aus und stimmte von neuem „Meßlieder“ an. Weiber werfen sich auf den Boden und wälzen sich wie wilde Thiere, der Priester schreitet durch die Menge hin und bespritzt sie mit „Weihwasser“. Dann hält er an die wildgewordene Menge eine „Predigt“, so schändlich, scheußlich, daß der Zeitungsschreiber mit Entsetzen die Hinterthür sucht: er hatte genug gesehen und gehört, er wollte nicht auch noch Augenzeuge sein der schamlosen Unzüchtigkeit, mit der die „schwarze Messe“ ihren Abschluß findet. Er bezeugt, daß die ganze Sache durchaus nicht ein Possenspiel auf die katholische Messe gewesen sei, noch etwa eine Verspottung des Teufels, sondern daß es allen Theilnehmern ein bitterer Ernst gewesen sei. — „Für dem Teufel uns bewahr!“ Fr. M.
